Indeed, we have already noted the inclusion @@M0@@. For the reverse inclusion, note that @@M1@@ belongs to @@M2@@ if and only if the sequence of @@M3@@ (@@M4@@) is quasi-compact\,; since @@M5@@ is quasi-compact, it suffices for this that the sequence of @@M6@@ have the same property\,; this gives the assertion.

\sgaStableTarget{sga6.e13.lem.4-3}\textbf{Lemma 4.3.} \emph{Let @@M7@@, @@M8@@ be proper over a scheme @@M9@@ and @@M10@@ a surjective morphism (resp. the inclusion morphism of the zero scheme of a section @@M11@@ of an ample sheaf @@M12@@, under the hypotheses of 3.8). Then}
@@M13@@

Indeed, the assertion is a formal consequence of 4.2, (XII 1.2), 3.1, and 3.5 (resp. 3.8).

\sgaStableTarget{sga6.e13.para.4-4}\textbf{4.4.} Let @@M14@@ be an algebraically closed field, @@M15@@ a proper @@M16@@-scheme, and @@M17@@ an invertible sheaf on @@M18@@. We say that @@M19@@ is \uline{algebraically equivalent} to @@M20@@ (resp. @@M21@@-\uline{equivalent} to @@M22@@) if there exist a connected @@M23@@-scheme @@M24@@ of finite type, an invertible sheaf @@M25@@ on @@M26@@, and two @@M27@@-points @@M28@@ of @@M29@@ such that @@M30@@ and @@M31@@ (resp. if there exists a nonzero integer @@M32@@ such that @@M33@@ is algebraically equivalent to @@M34@@), or, what is easily seen to be equivalent, if the @@M35@@-point @@M36@@ of @@M37@@ corresponding to @@M38@@ belongs to @@M39@@ (resp. to @@M40@@).

Recall also that @@M41@@ is said to be \uline{numerically equivalent} to @@M42@@ if, for every integral closed curve @@M43@@ in @@M44@@, the intersection number @@M45@@ is @@M46@@ (see 2.1 or X 4.5).

